\section{Technical Background}
\label{sec:background}
\noindent \textbf{Radar Basics}
\label{sec:background-mm}
\begin{figure}[t]
  \centering
   \includegraphics[width=\linewidth]{figures/diffusion-fmcw-1.pdf}
   \caption{(a) Frequency Modulated Continuous Waveform, (b) the reflected signal is a delayed version of the transmitted signal, to resolve distance. (c) Specular reflections dominate radar wave propagation vs (d) diffused reflections that dominate vision, (e) the radar reflections modeled to incorporate specular reflections and signal spread.}
   \vspace{-1.5em}
   \label{fig:fmcw}
\end{figure}
A mmWave radar works by transmitting a wireless
signal and receiving back reflections that come from reflectors in the scene.
It operates in the millimeter-wavelength frequency bands, and uses Frequency Modulated Continuous Wave (FMCW) and antenna arrays to help resolve spatial ambiguity.
As shown in Fig.~\ref{fig:fmcw}(a), the transmitted FMCW is a function that is linearly increasing in frequency over time.
Considering a single reflector and a single antenna, the received chirp is simply a time-delayed version of the transmitted chirp, see Fig.~\ref{fig:fmcw}(b).
To resolve range ambiguity, the received chirp is multiplied with the conjugate of the transmitted chirp and can be expressed as a complex function:
\begin{equation}
\label{eq:fmcw}
s(t) =  A \cdot e^{-j 2 \pi(f+kt)   d /c}  = A \cdot e^{-(j 2\pi k \tau )t} \cdot e^{-j 2\pi f \tau}
\end{equation}
where $A$ is the signal amplitude, $d$ is the round-trip propagation distance, $c$ is the speed of light, $\tau=d/c$ is the round-trip delay, $f$ is the starting frequency, and $k$ is the chirp slope. For multiple reflectors, we simply receive the linear combination of all the reflections.

\begin{figure}[t]
    \centering
    \begin{subcaptionbox}{Lens-based vs lensless imaging \label{fig:lensless}}[0.4\columnwidth]
        {\includegraphics[width=\linewidth]{figures/lens-less.pdf}}
    \end{subcaptionbox}
    \hfill
    \begin{subcaptionbox}{Example MF image from measured radar data.\label{fig:heatmap}}[0.55\columnwidth]
        {\includegraphics[width=\linewidth]{figures/mf_example.pdf}}
    \end{subcaptionbox}
    \caption{a) Comparison between lens-based imaging and lensless imaging models.
b) Matched filter image of the metal bunny, summed along the vertical, horizontal, and depth axis.
}
\vspace{-1.5em}
    \label{fig:mf-example}
\end{figure}

\vspace{-0.1em}
\noindent \textbf{Reflections}
Unlike light, whose wavelength (order of \( 10^{-7} \) meters) is much smaller than surface irregularities, resulting in predominantly diffused reflection, RF signals have much longer wavelengths (order of \( 10^{-3} \) meters).
In this case, the majority of surfaces appear relatively smooth to the RF signal, and specular reflections becomes more prominent than diffused scattering.
For specular reflections, the angle of incidence is the same as the angle of reflection, while diffused reflections tend to scatter at different angles. The difference is illustrated in Fig.~\ref{fig:fmcw}(c,d).

In this paper, we primarily expect specular reflections.
We follow a shifted Lambertian model~\cite{richards2010principles, ren2020diffuse}, where the power of the reflections that return to the receiver, is scaled according to the normal and the signals reflection angle, see Fig.~\ref{fig:fmcw}(e).
The received power is modulated by the directional scaling factor \( (\omega_o \cdot \omega_r) \), where \( (\omega_o\) is the ideal specular vector and \( \omega_r \) is the vector from the reflecting surface to the receiver (see Appendix~\ref{sec:app-radar-ref}).
The reflection strength also depends on material properties and the thickness of the object, which we model using a reflective coefficient \( a \).

An additional power decay factor of \({(4\pi u)^2}\) is divided from the received reflection amplitude (\({(4\pi u)^4}\) for power), where \(u\) is the distance between the point and antennas, to accurately model the free space path loss (see Appendix~\ref{sec:app-radar-pwr}). Given an input signal with amplitude $ A_{\text{tx}} $, the received signal amplitude $ A_{\text{rx}} $ can be expressed as:
\begin{equation}
\label{eq:reflection_angle}
A_{\text{rx}} \propto \frac{a}{(4\pi u)^2} A_{\text{tx}}\, (\omega_o \cdot \omega_r)
\end{equation}

\vspace{-0.1em}
\noindent \textbf{Lensless imaging}
Unlike traditional visual imaging systems that rely on lenses, radar imaging is inherently a \textbf{lensless} imaging technique.
As illustrated in Fig.~\ref{fig:lensless}, a pinhole in visual systems acts as a physical filter, allowing only one ray per scene point to reach the CMOS sensor while blocking unrelated rays.
In photography, increasing the aperture size allows more light rays to enter, but this also leads to blur and noise in the image.
As the aperture approaches an infinite size, the system transitions into a lensless imaging model.
Lensless imaging, such as with radar antenna arrays, does not employ any physical pinhole—\textit{all} incoming rays are received by all antennas.

In contrast, radar imaging relies on algorithmic processing to reconstruct the scene.
The core idea involves applying a digital filter to remove unrealistic signals and keep the most relevant ones, leveraging the time-varying and phase-delay information of the received signals.
Each received signal, arriving via a different path, carries a distinct phase shift defined by its propagation delay.
The algorithm used for this process is called the \textbf{matched filter}.

\vspace{-0.1em}
\noindent \textbf{Differentiable Matched Filter}
This algorithm is implemented by correlating the received signal with an ideal reference signal.
From Eq.~\ref{eq:fmcw}, the reference signal for a single reflector and a single antenna is known.
For multiple antenna measurements, we have:
\begin{equation}
\label{eq:mf-sm}
P(\mathbf{x}) = \left| \left| \sum_{i = 1}^{N_{\text{ant}}} \sum_{t} s(i, t) \, e^{j 2\pi k \tau_i t} \, e^{j 2\pi f \tau_i} \right| \right|,
\end{equation}
where \( N_{\text{ant}} \) is the number of antennas, \( \tau_i \) represent the round-trip delay corresponding to the \( i \)-th antenna, and \( s(i, t) \) denotes the received signal at time \( t \) from the \( i \)-th antenna. We omit the transmitted amplitude \( A_{\text{tx}} \) since it is constant.
This process effectively averages out signals that are inconsistent with the hypothesized signal model, a matched filter (MF) power image is shown in Fig.~\ref{fig:heatmap}.

To enable integration into a differentiable rendering framework, we also compute the backpropagation of MF.
Both the forward and backward passes are implemented in parallel on the GPU.
All implementation details are provided in Appendix~\ref{sec:app-mf} and \ref{sec:app-alg}.
